\documentclass[12pt,aspectratio=169,t]{beamer}
\usetheme[numbering=none,progressbar=none]{metropolis}
\usepackage{booktabs}
\usepackage[compatibility=false]{caption}

\defaultfontfeatures{Ligatures=TeX}

\usepackage{unicode-math}
\setmathfont{Latin Modern Math}

% High-contrast text colors shared with the manuscript link palette.
\definecolor{deepteal}{RGB}{40,72,82}
\definecolor{accentred}{RGB}{123,0,0}
\setbeamertemplate{headline}{}
\setbeamertemplate{footline}{}
\setbeamerfont{frametitle}{size=\Large}
\setbeamerfont{section title}{size=\Huge, series=\bfseries}
\setbeamertemplate{frametitle}{
  \vspace{0.5cm}
  \insertframetitle
  \vspace{0.3cm}
}
\addtobeamertemplate{frametitle}{}{\vspace{-0.5em}}
\setbeamertemplate{section page}{
  \centering
  \vfill
  {\usebeamerfont{section title}\usebeamercolor[fg]{section title}\insertsection\par}
  \vfill
}
\setbeamercolor{background canvas}{bg=white}
\setbeamercolor{normal text}{fg=black}
\setbeamercolor{frametitle}{fg=black, bg=white}
\setbeamercolor{title}{fg=black}
\setbeamercolor{subtitle}{fg=black}
\setbeamercolor{author}{fg=black}
\setbeamercolor{date}{fg=black}
\setbeamercolor{progress bar}{fg=deepteal, bg=white}
\setbeamercolor{alerted text}{fg=accentred}
\setbeamerfont{alerted text}{series=\bfseries}
\setbeamercolor{example text}{fg=deepteal}
\setbeamercolor{section title}{fg=accentred}


\title{Broad Validity of the First-Order Approach in Moral Hazard}
\author{Eduardo M. Azevedo and Ilan Wolff (Wharton)}
\date{March 2, 2026}

\begin{document}

\begin{frame}[plain]
  \titlepage
\end{frame}

\begin{frame}{The Moral Hazard Problem (1970s)}
  \begin{itemize}
    \item Principal hires agent. Output
    \[
      y \sim f(y | a)
    \]
    \item Agent utility
    \[
      u(\text{payment}) - c(a)
    \]
    \item Principal's problem is to find optimal contract \textcolor{accentred}{$w(y)$}.
  \end{itemize}
\end{frame}

\begin{frame}{What is known}
  \begin{itemize}
    \item The moral hazard problem has a simple solution if the first-order approach is valid \textcolor{deepteal}{(Holmstrom 78)}.
    \item Unfortunately, FOA is only valid under restrictive assumptions \textcolor{deepteal}{(Mirrlees 75, Rogerson 85, Jewitt 88, Sinclair-Desgagné 94)}.
    \item So theoretical optimal contracts are \textcolor{accentred}{intractable}, \textcolor{accentred}{complex}, and \textcolor{accentred}{not realistic}.
    \item Motivates huge literature proving linearity. Theoretical literature often assumes FOA anyway. Applied literature often assumes two possible outputs or linearity.
  \end{itemize}
\end{frame}

\begin{frame}
  \begin{columns}[T,onlytextwidth]
    \column{0.40\textwidth}
    \textbf{Simplest example}
    \begin{itemize}
      \item Gaussian output: $y \sim \mathcal N(a, \sigma^2)$, $\sigma = \$50{,}000$.
      \item Log utility, initial wealth $\$50{,}000$.
      \item Limited liability.
      \item Quadratic cost.
    \end{itemize}

    \textbf{Results at odds with received wisdom.}

    \column{0.56\textwidth}
    \centering
    \includegraphics[height=0.95\textheight]{../figures/log-gaussian-sigma=50.0/pp_stacked.pdf}
  \end{columns}
\end{frame}

\begin{frame}{What is actually true}
  \begin{itemize}
    \item The FOA holds except for very low reservation utility.
    \item Key theorem assumptions: at least moderate risk aversion and $f_{aa}(y|a) \geq 0$ in a large enough region of low scores.
    \item Side effects:
    \begin{itemize}
      \item \textcolor{accentred}{Tractable} (analytically and computationally trivial).
      \item \textcolor{accentred}{Simple} (intuitive formula, linear w/ log utility and exponential family with linear sufficient statistic).
      \item \textcolor{accentred}{Realistic}.
    \end{itemize}
    \item Interactive examples: \url{https://eduardomazevedo.github.io/azevedo-wolff-2025-site/}.
  \end{itemize}
\end{frame}

\begin{frame}{Literature is mathematically correct}
  \begin{itemize}
    \item There is not a mathematical error in key prior results.
    \item The literature is correct that any sufficient condition that guarantees that FOA holds for any reservation utility is extremely restrictive. Gaussian example above would be ruled out because FOA fails at monopsony reservation utility.
    \item But sufficient conditions are mild if we include sufficiently high reservation utility. See all the examples above and theorem below.
  \end{itemize}
\end{frame}

\begin{frame}{Roadmap}
  \begin{itemize}
    \item \textcolor{alert}{Model}
    \item Main result
    \item Applications
  \end{itemize}
\end{frame}

\begin{frame}{Model}
  Agent has initial wealth $w_0$, utility $u(w_0+w)-c(a)$ from payment $w\geq 0$ and action $a\in\mathcal A$. Reservation utility $\bar U$.

  Output $y$ has density $f(y\mid a)$.

  Contract is $w\colon\mathcal Y\to\mathbb R_+$. Feasible contracts: $\mathcal W$.

  \begin{align}
    U(w, a) &= \int u\bigl(w_0+w(y)\bigr) \cdot f(y\mid a)\,\mathrm{d}y - c(a) \\
    C(w, a) &= \int w(y) \cdot f(y\mid a)\,\mathrm{d}y
  \end{align}
\end{frame}

\begin{frame}{Definitions: FOA}
  \textbf{GIC} (global incentive compatibility): choosing $a_0$ is optimal for the agent.

  \textbf{LIC} (local incentive compatibility): $U_a(w, a_0) = 0$.

  \textbf{Cost minimization problem:} minimize expected wage subject to IR and GIC.

  \textbf{Relaxed cost minimization problem:} minimize expected wage subject to IR and LIC.

  \textbf{FOA holds for $(a_0, \bar{U})$} if the solution to the relaxed cmp solves the cmp.
\end{frame}

\begin{frame}{Definitions}
  \begin{itemize}
    \item Score $s(y) := f_a(y\mid a_0)/f(y\mid a_0)$.
    \item Sufficient risk aversion: with absolute risk aversion $A$ and absolute prudence $P$,
    \[
      \frac{A(W)}{P(W)}\geq\frac{1}{3}
      \qquad\text{and}\qquad
      \lim_{W\to\infty}\frac{A(W)}{u'(W)}=\infty.
    \]
    Includes CARA and CRRA with $\gamma>1$.
  \end{itemize}
\end{frame}

\begin{frame}{Safe capacity and curvature}
  A cutoff $q\leq 0$ is \textbf{safe} if $f_{aa}\geq 0$ on $\{s\leq q\}$ for every $a$.
  Let $\mathcal Q_{\mathrm{safe}}$ be the set of safe cutoffs, and $\bar q_{\mathrm{safe}}:=\sup\mathcal Q_{\mathrm{safe}}$.
  \[
    \mathrm{Cap}_{\mathrm{safe}}
    :=-
    \int_{\{s(y)<\bar q_{\mathrm{safe}}\}}
    s(y)f(y\mid a_0)\,d\nu(y).
  \]
  \[
    \mathrm{Curv}_{\mathrm{safe}}
    :=\inf_{q\in\mathcal Q_{\mathrm{safe}}}
      \sup_{a\in\mathcal A}
      \int_{\{s(y)>q\}}
      \bigl[s(y)f_{aa}(y\mid a)\bigr]_+\,d\nu(y).
  \]
\end{frame}

\begin{frame}{Outline}
  \begin{itemize}
    \item Model
    \item \textcolor{alert}{Main result}
    \item Applications
  \end{itemize}
\end{frame}

\begin{frame}{Theorem 1}
  Suppose either:
  \begin{enumerate}
    \item[(i)] sufficient risk aversion, $\underline c''>0$, $\mathrm{Curv}_{\mathrm{safe}}<\infty$, and
    \[
      \bigl[\lim_{W\to\infty}u(W)-u(w_0)\bigr]\mathrm{Cap}_{\mathrm{safe}}>c'(a_0);
    \]
    \item[(ii)] log utility and $\mathrm{Curv}_{\mathrm{safe}}/\mathrm{Width}_{\mathrm{safe}}<\underline c''$.
  \end{enumerate}
  Then there exists $U^*<U_R$ such that, for every $\bar U\in[U^*,U_R)$, the first-order approach is valid.
\end{frame}

\begin{frame}{Proof part 1: Relaxed problem}
  Write Lagrangian:
  \[
    \mathcal{L}(w, \lambda, \mu) = C(w,a_0) + \lambda (\bar{U} - U(w,a_0)) - \mu U_a(w,a_0)
  \]
  \begin{align}
    C(w,a_0)   &= \int w(y) \cdot f\phantom{_a}(y\mid a_0)\,\mathrm{d}y \\
    U(w,a_0)   &= \int u\bigl(w_0+w(y)\bigr) \cdot f\phantom{_a}(y\mid a_0)\,\mathrm{d}y - c(a_0) \\
    U_a(w,a_0) &= \int u\bigl(w_0+w(y)\bigr) \cdot f_a(y\mid a_0)\,\mathrm{d}y - c'(a_0)
  \end{align}
\end{frame}

\begin{frame}{Proof part 1: Relaxed problem}
  FOC:
  \[
    \frac{1}{u'\bigl(w_0+w(y)\bigr)} = \lambda + \mu s(y)
  \]
  Solve for $w(y)$:
  \[
    w(y) = L\bigl(\lambda + \mu s(y)\bigr) - w_0,
  \]
  where $m_0:=1/u'(w_0)$, $L_0(z):=(u')^{-1}(1/z)$, and
  \[
    L(z) := L_0(z\vee m_0).
  \]
\end{frame}

\begin{frame}{Proposition 1 ($\approx$ Holmstrom 78)}
  The relaxed problem has an almost everywhere unique solution $w^*_{\bar U}$. There exist $\lambda$, $\mu$ such that
  \[
    w^*_{\bar U}(y) = L\bigl(\lambda + \mu s(y)\bigr) - w_0.
  \]
\end{frame}

\begin{frame}{Proof part 2: concavity}
  \textbf{Proposition 2:}

  Under either set of conditions in Theorem~1, there exists $\lambda_0<\infty$ such that every locally incentive-compatible canonical contract with $\lambda\geq\lambda_0$ obeys
  \[
    U_{aa}\bigl(w_{\lambda,\mu}, a\bigr) \leq 0
    \qquad\text{for every }a\in\mathcal A.
  \]
\end{frame}

\begin{frame}
  \begin{columns}[T,onlytextwidth]
    \column{0.40\textwidth}
    \textbf{Key idea:} kink moves to the left when $\bar{U}$ goes up.

    \column{0.56\textwidth}
    \centering
    \includegraphics[height=0.95\textheight]{../figures/log-gaussian-sigma=50.0/cm_stacked.pdf}
  \end{columns}
\end{frame}

\begin{frame}{Heuristic derivation}
  \[
    c'(a_0)
    =\int\bigl[u\bigl(L(\lambda+\mu s(y))\bigr)
                       -u\bigl(L(\lambda)\bigr)\bigr]s(y)
      f(y\mid a_0)\,\mathrm{d}y.
  \]
  If the kink is to the right of a cutoff $q_0<0$, then
  \[
    c'(a_0)
    \geq \bigl[u\bigl(L(\lambda)\bigr)-u(w_0)\bigr]
      \Bigl(-\!\int_{\{s\leq q_0\}}s\,f\,\mathrm{d}y\Bigr).
  \]
  As $\lambda\to\infty$, the right-hand side exceeds $c'(a_0)$ for a safe cutoff, so the kink must move into the safe low-score region.
\end{frame}

\begin{frame}{Heuristic derivation}
  Centering $U_{aa}$ and dropping the safe region (where the integrand is $\leq 0$),
  \[
    U_{aa}
    \leq\int_{\{s>q_1\}}
      \bigl[u\bigl(L(\lambda+\mu s)\bigr)-u\bigl(L(\lambda)\bigr)\bigr]
      f_{aa}\,\mathrm{d}y
      -c''(a).
  \]
  The kink-left-of-$q_0$ bound keeps $\mu$ from growing too fast relative to $\lambda$. Utility conditions tame sensitivity above $q_1$, and $\mathrm{Curv}_{\mathrm{safe}}$ bounds positive distributional curvature. Remaining curvature is eventually dominated by $\underline c''$.
\end{frame}

\begin{frame}{Heuristic derivation}
  With utility link $\ell(z):=u(L(z))$ and $L:=\lim_{z\to\infty}z\ell_0'(z)$,
  \begin{eqnarray*}
    U_{aa}
    &\leq&
    \displaystyle\int_{\{s>q_1\}}
      \bigl[\ell(\lambda+\widetilde\mu s)-\ell(\lambda)\bigr]
      f_{aa}\,\mathrm{d}y
      -c''(a) \\
    &\leq&
    \widetilde\mu\,\ell_0'(\lambda+\widetilde\mu q_1)\,\mathrm{Curv}(q_1)-\underline c'' \\
    &\leq&
    \mu_{q_0}\,\ell_0'(\lambda+\mu_{q_0}q_1)\,\mathrm{Curv}(q_1)-\underline c'' \\
    &=&
    \frac{\mu_{q_0}}{m_0+\mu_{q_0}(q_1-q_0)}
    \bigl[(m_0+\mu_{q_0}(q_1-q_0))\ell_0'(m_0+\mu_{q_0}(q_1-q_0))\bigr]
    \mathrm{Curv}(q_1)-\underline c'' \\
    &\to&
    \dfrac{L\,\mathrm{Curv}(q_1)}{q_1-q_0}-\underline c''.
  \end{eqnarray*}
\end{frame}

\begin{frame}{Outline}
  \begin{itemize}
    \item Model
    \item Main result
    \item \textcolor{alert}{Applications}
  \end{itemize}
\end{frame}

\begin{frame}{Calculus}
  \begin{itemize}
    \item Optimal contracts have simple and intuitive closed-form solutions.
    \item Calculus of moral hazard has two parts: wage link function and score.
  \end{itemize}
\end{frame}

\begin{frame}[plain]
  \centering
  \begin{table}[ht]
    \centering
    \caption{Utility Functions and Wage Link Functions}
    \label{tab:utility_functions}
    \renewcommand{\arraystretch}{1.5}
    \setlength{\tabcolsep}{10pt}
    \begin{tabular*}{\textwidth}{@{\extracolsep{\fill}}lcc}
        \toprule
        & \multicolumn{1}{c}{Utility Function}
        & \multicolumn{1}{c}{Wage Link Function} \\
        & \(u(W)\) & \(L(z)\) \\
        \midrule
        Log
        & \(\log W\)
        & \(z\vee w_0\) \\
        CRRA
        & \(\dfrac{W^{1-\gamma}}{1-\gamma}\)
        & \(\bigl(z\vee w_0^\gamma\bigr)^{1/\gamma}\) \\
        CARA
        & \(-\dfrac{\exp(-\alpha W)}{\alpha}\)
        & \(\dfrac{\log\bigl(z\vee\exp(\alpha w_0)\bigr)}{\alpha}\) \\
        \bottomrule
    \end{tabular*}
    \captionsetup{font=footnotesize}
    \caption*{\textit{Note:} Utility $u$ is defined over total wealth $W$.  The
    wage link function $L$ gives canonical total wealth; subtracting initial wealth gives
    the associated wage,
    $w_{\lambda,\mu}(y)=L(\lambda+\mu s(y))-w_0$.  For CRRA,
    $\gamma>0$ and $\gamma\neq1$; for CARA, $\alpha>0$.}
\end{table}

\end{frame}

\begin{frame}[plain]
  \centering
  \begin{minipage}{0.90\textwidth}
    \begingroup
    \let\small\scriptsize
    \let\footnotesize\tiny
    \input{../tex/tables/distributions}
    \endgroup
  \end{minipage}
\end{frame}

\begin{frame}
  \begin{columns}[T,onlytextwidth]
    \column{0.40\textwidth}
    \textbf{Piecewise linear contracts}
    \medskip

    Optimal contract is piecewise linear for log utility + exponential family with linear sufficient statistic.
    \medskip

    Prior results: \textcolor{deepteal}{Opp 25} (independently), earlier papers without existence \textcolor{deepteal}{(Hemmer 99, Christensen 05)}.
    \medskip

    Global linearity in some cases.

    \column{0.56\textwidth}
    \centering
    \includegraphics[height=0.95\textheight]{../figures/log-gaussian-sigma=50.0/pp_stacked.pdf}
  \end{columns}
\end{frame}

\begin{frame}
  \begin{columns}[T,onlytextwidth]
    \column{0.40\textwidth}
    \textbf{Piecewise linear contracts}
    \medskip

    Optimal contract is piecewise linear for log utility + exponential family with linear sufficient statistic.
    \medskip

    Prior results: \textcolor{deepteal}{Opp 25} (independently), earlier papers without existence \textcolor{deepteal}{(Hemmer 99, Christensen 05)}.
    \medskip

    Global linearity in some cases.

    \column{0.56\textwidth}
    \centering
    \includegraphics[height=0.95\textheight]{../figures/log-poisson/pp_stacked.pdf}
  \end{columns}
\end{frame}

\begin{frame}{Numerical methods}
  \begin{itemize}
    \item Results suggest efficient algorithms.
    \item Consider a finite set of GIC constraints at $\boldsymbol{\hat a} = (\hat a_1,\ldots,\hat a_n)$.
    \item With multipliers $\theta := (\lambda,\mu,\boldsymbol{\hat\mu})$, the optimal contract is
    \[
      w(y | \theta, \hat a) = L\left(\lambda + \mu s(y) + \sum_i \hat\mu_i\left(1-\frac{f(y|\hat a_i)}{f(y|a_0)}\right)\right)-w_0.
    \]
    \item Lagrange dual $\mathcal D(\theta, \hat a)$ is fast to compute if we cache these terms.
    \item Analytic gradient from Danskin's envelope theorem.
  \end{itemize}
\end{frame}

\begin{frame}{Numerical methods: algorithm}
  \begin{enumerate}
    \item Start with $\hat a = \emptyset$.
    \item Maximize dual to solve relaxed problem, and add best deviation action to $\hat a$ if needed.
    \item Repeat until no deviation is found.
    \item In case of trouble, fall back to convex optimizer to find more deviations for $\hat a$.
  \end{enumerate}
\end{frame}

\begin{frame}
  \begin{columns}[T,onlytextwidth]
    \column{0.40\textwidth}
    \textbf{Performance}
    \medskip

    Active-set algorithm runs in about \textcolor{accentred}{3--14 ms (CMP)} and \textcolor{accentred}{232--240 ms (PP)} vs. hours reported in the literature.
    \medskip

    Why it is fast:
    \begin{itemize}
      \item Often terminates in one iteration.
      \item Dual calls use analytic formulas + caching.
      \item Low-dimensional dual with analytic gradients + warm starts.
    \end{itemize}

    \column{0.56\textwidth}
    \centering
    \includegraphics[height=0.95\textheight]{../figures/solver_comparison/stacked_comparison.pdf}
  \end{columns}
\end{frame}

\begin{frame}
  \begin{columns}[T,onlytextwidth]
    \column{0.40\textwidth}
    \textbf{Key assumption:} existence of a safe score region.
    \medskip

    Example: $t$-distributed outcome, tail coefficient $\nu = 1.15$.
    \medskip

    No nonempty safe low-score region, and FOA often fails even for high reservation utility.

    \column{0.56\textwidth}
    \centering
    \includegraphics[height=0.95\textheight]{../figures/log-t-sigma=20.0/pp_stacked.pdf}
  \end{columns}
\end{frame}

\begin{frame}{Conclusion}
  \begin{itemize}
    \item FOA is usually valid, we just have to assume high reservation utility.
    \item Tractable, intuitive, realistic solutions---sometimes linear contracts.
  \end{itemize}
  \vspace{4em}
  \centering
  \textcolor{deepteal}{\Huge Thank you!}
\end{frame}

\end{document}
